\section{Introduction}




Multi-armed bandits (MABs) have garnered significant attention as they provide a well-defined framework and techniques for investigating the exploration-exploitation trade-off in sequential decision-making problems. In the context of MAB problems, a decision-maker sequentially selects actions from a given set and observes corresponding rewards. This paper specifically focuses on the stochastic linear bandit (LB) model, which encompasses MAB and contextual bandit problems as special cases. In this model, the system parameter $\param$ represents an unknown vector in $\RR^d$, while the actions, also vectors in $\RR^d$, yield rewards equal to the inner product of $\param$ and the chosen action, along with mean-zero, subgaussian noise.
	The objective of a policy is to maximize the cumulative reward based on the observed data up to the decision time. The policy's performance is measured by the cumulative regret, which quantifies the difference between the total rewards achieved by the policy and the maximum achievable reward. Achieving this objective necessitates striking a balance between exploration and exploitation. In the context of LB, this entails selecting actions that aid in estimating the true parameter $\param$ accurately while obtaining optimal rewards. Various algorithms based on the Optimism in the Face of Uncertainty (OFU) principle have been developed to address this challenge, wherein the optimal action is chosen based on the upper confidence bound (UCB) \citep{lai1985asymptotically}. Other popular strategies include Thompson sampling (TS), a Bayesian heuristic introduced by \cite{thompson1933likelihood} that employs randomization to select actions according to the posterior distribution of reward functions. Additionally, the Greedy policy that selects the myopically best action  is shown to be effective in contextual bandits \citep{kannan2018smoothed,raghavan2018externalities,hao2020adaptive, bastani2021mostly}.

In the linear bandit setting, two types of regret are considered. The Bayesian regret is applicable when the parameter $\param$ is random, and the regret is averaged over three sources of randomness: the observation noise, the randomized algorithm, and the random parameter $\param$. \cite{russo2014learning} and \cite{dong2018information} establish an $\tilde{\cO}(d \sqrt{T})$ upper bound for the Bayesian regret of the TS heuristic referred to as LinTS, which matches the minimax optimal bound shown by \cite{dani2008stochastic}, where $\tilde{\cO}$ refers to asymptotic order, up to logarithmic factors. On the other hand, the frequentist regret corresponds to scenarios where the underlying parameter $\param$ is fixed, and the expectation is taken only with respect to the randomness of noise and algorithm. % Estimating the frequentist regret is more challenging as it requires considering all unfavorable cases.

In the frequentist setting, the optimism-based algorithm proposed by \cite{abbasi2011improved}, known as Optimism in the Face of Uncertainty Linear Bandit Algorithm (OFUL), achieves a $\tilde{\cO}(d\sqrt{T})$ frequentist regret bound, which matches the minimax optimal Bayesian bound. However, for a frequentist variant of LinTS referred to as TS-Freq, \cite{agrawal2013thompson} and \cite{abeille2017linear} provide a $\tilde{\cO}(d\sqrt{dT})$ upper bound for the frequentist regret by inflating the posterior distribution. This regret bound falls short of the optimal rate by a factor of $\sqrt{d}$. Recent work by \citet{hamidi2020frequentist} confirms that this inflation is necessary, and thus, the frequentist regret of LinTS cannot be improved. While no general theoretical guarantee exists for the Greedy algorithm in linear bandit problems, it is known that both Greedy and LinTS perform suboptimally in the asymptotic regime \citep{lattimore2017end}.

Despite the theoretical gaps in performance, LinTS has shown strong empirical performance, indicating that the inflation of the posterior distribution may be unnecessary in most problem instances, and unfavorable instances are unlikely to occur frequently. Similarly, the Greedy algorithm works well in typical instances 
\citep{bietti2021contextual}. Additionally, while OFU-based algorithms are computationally expensive (generally NP-hard as discussed by \cite{dani2008stochastic,russo2014learning,agrawal2019recent}), LinTS and Greedy are known for their computational efficiency. This unsettling disparity between theoretical, computational, and empirical performance has motivated our investigation into the following two questions:  Is it possible to identify problematic instances where LinTS and Greedy could potentially fail, and apply a ``course-correction'' to ensure competitive frequentist regret bounds? And can this be achieved without compromising their impressive empirical performance and computational efficiency?
%
In this paper, we provide positive answers to both questions. Specifically, we make the following \emph{contributions}. 

    \begin{itemize}
    \item Our technique conducts a novel real-time geometric analysis of the $d$-dimensional confidence ellipsoid. This technique is essential in maximizing the utilization of historical data, thereby enabling the ``course-correction''.
    %\item We propose a family of algorithms called \emph{POFUL} (including OFUL, LinTS, TS-Freq, and Greedy as special cases) and derive a general instance-dependent frequentist regret bound for them.
\item We derive a general instance-dependent frequentist regret bound for the proposed \emph{POFUL} algorithms, a family of algorithms including OFUL, LinTS, TS-Freq, and Greedy as special cases. The bound can be efficiently calculated in a data-driven manner based on the observation history.
\item  We propose variants of LinTS and Greedy that enjoy the minimax optimal frequentist regret, while retaining most of the desirable properties such as computational efficiency.
    \end{itemize}

%\begin{itemize}
%	\item We propose a generalized version of  OFUL that we call POFUL (Pivot OFUL), which includes OFUL, LinTS, TS-Freq, and Greedy as special cases.
%	\item We derive an instance-dependent frequentist regret bound for POFUL that can be efficiently calculated in a data-driven manner based on the observation history.
%	\item We propose variants of LinTS and Greedy that enjoy matching frequentist regret as the Bayesian minimax regret of order $d\sqrt{T}$ up to logarithmic factors. 
%\end{itemize}


The distinguishing feature of our study from the existing literature is the geometric analysis of the $d$-dimensional confidence ellipsoid. 	This is in stark contrast to techniques outlined in previous studies that may inadvertently discard valuable information as a consequence of their overly narrow concentration on comparatively rudimentary summary statistics.

\paragraph{Related literature.} The ideas presented in this paper bear similarities to existing literature that explores different forms of exploration \citep{russo2016information,kirschner2018heteroscedastic} or recent works by \citep{bastani2021mostly,pacchiano2020regret,hamidi2020general} that aim to utilize observed data for reducing exploration. However, the methodology and the manner in which the data is utilized in this paper differ significantly. For instance, \cite{bastani2021mostly} examines the minimum eigenvalue of the sample covariance matrix which is a single-parameter summary statistic of the observed data, while 
 \cite{hamidi2020general} incorporates information from the (one-dimensional) confidence intervals around the reward values. In contrast, our approach leverages the complete geometric information of the $d$-dimensional confidence ellipsoid.


  
  
  \begin{comment}
\begin{itemize}
	\item We derive an instance-dependent frequentist regret bound for TS (without inflation), which can be explicitly calculated along the execution path of the algorithm. 
	\item We introduce a new algorithm (TS-MR) that mimics TS, and only deviates from TS when it detects TS is incurring a large regret. And we prove that TS-MA only suffers a constant factor in the frequentist regret.
	\item We show in simulations that our regret bounds are tight in most cases, providing insights into TS's empirical success. Also, we show in simulations that TS-MA outperforms TS-Freq in terms of both regret and empirical performance.
	\item Our techniques study the geometric property of the posterior distribution of TS, which is a new perspective to understand the performance of TS and can be of independent interest.


\end{itemize}
\end{comment}
